% Generated from the canonical article appendix; do not edit.

This appendix proves the seed and query assertions in Chapter~\ref{comp:D} in three purely rational stages:
isolating the algebraic root $(c,d)$, verifying symbolic closure of
the three-step matrix and iterate cycle, and enclosing all nine trial queries in
exact rational interval arithmetic. All decimals below denote exact terminating
rationals.

\section{Existence and uniqueness of the defining root}
We choose the third gradient $g_2=(c,d)^T$ so that the last two accepted step
lengths in the symbolic cycle of the next section satisfy $\beta(c,d)=1/4$ and
$\gamma(c,d)=1/16$. Clearing positive denominators yields the polynomial system
\begin{equation}\label{comp:seed_polynomials}
 P(c,d)=
 \begin{pmatrix}264192RK-UV(U+129)\\16VTK-RUY\end{pmatrix}=0,
\end{equation}
where
\begin{align*}
 U&=127c+16d,&V&=16c-127d-16,&R&=c-8d,\\
 K&=8c^2+8d^2+d-8,&T&=c^2-c+d^2,&Y&=(c-1)^2+d^2.
\end{align*}
We isolate the root inside the closed $\ell_\infty$ box
$\|(c,d)-\bar r\|_\infty\le10^{-30}$ centered at
\begin{align*}
 \bar r_1&=0.31848142762711863064645131250290703717912205023605,\\
 \bar r_2&=-1.02523211967574872548135087001300414284511909318444.
\end{align*}

\begin{lemma}[Rational root certificate]\label{comp:root_certificate}
The closed box $\|(c,d)-\bar r\|_\infty\le10^{-30}$ contains a unique
root of $P$. Throughout this box all of $U,V,R,K,T,Y,U+129$ are strictly
positive, and~\eqref{D:rho} holds.
\end{lemma}
\begin{proof}
Set $D=DP(\bar r)$ and let $A=\operatorname{adj}(D)/\det D$ be its $2\times2$
rational inverse. Direct rational evaluation of the eighteen Taylor monomials of
$P$ of total degrees $2\le|\nu|\le5$ at $\bar r$ certifies the three exact
inequalities
\begin{equation}\label{comp:root_arithmetic}
 \begin{gathered}
 \det D\ne0,\qquad
 \max_i |(AP(\bar r))_i|<10^{-50},\\
 \max_i\sum_{2\le|\nu|\le5}|a_{i,\nu}|\,|\nu|\le7000,
 \qquad
 a_{i,\nu}=-\frac{1}{\nu_1!\nu_2!}\sum_{j=1}^2A_{ij}\,\partial^\nu P_j(\bar r).
 \end{gathered}
\end{equation}
Because $AD=I$, the preconditioned Newton map $N(h)=h-AP(\bar r+h)$ has zero
linear part and exact polynomial expansion
$N_i(h)=-(AP(\bar r))_i+\sum_{2\le|\nu|\le5}a_{i,\nu}h^\nu$. On the closed ball
$\|h\|_\infty,\|k\|_\infty\le r=10^{-30}$, telescoping each monomial gives
$|h^\nu-k^\nu|\le|\nu|r\|h-k\|_\infty$ ($|\nu|\ge2$), so $N$ is a contraction
with Lipschitz constant at most $7000r\ll1$ and maps the ball into itself since
$\|N(h)\|_\infty\le10^{-50}+7000r^2<r$ (cf.\ Rump~\cite[Theorem~13.3]{NEG_RUMP10}).
Invertibility of $A$ identifies its unique fixed point with the unique root of
$P$ in the box.

Since the box lies inside $318/1000<c<319/1000$ and
$-1026/1000<d<-1025/1000$, interval evaluation gives strict positivity of
$U,V,R,K,T,Y,U+129$; for instance,
\[
 \frac{2397}{100}\le U\le\frac{24113}{1000},\quad
 \frac{119263}{1000}\le V\le\frac{59703}{500},\quad
 \frac{4259}{500}\le R\le\frac{8527}{1000}.
\]
Substituting these positive bounds into $\rho=129R/[V(U+129)]$ yields
$60/1000<\rho<61/1000$. In particular, all divisions used below are valid. The same interval rules on the root box give the thirty digits in \eqref{D:rho}.
\end{proof}

\section{Symbolic closure of the three-step matrix cycle}
\begin{proof}[Proof of the algebraic part of the seed and query assertions in Chapter~\ref{comp:D}]
From $s_0=(-1,0)^T$ and $g_1-g_0=(-2,-1/8)^T$, the first update is immediate:
$\mathcal B(I,g_0,g_1,1)=M_1$. For any $(c,d)$ where the factors
above are positive, define the rational step lengths and contraction ratio
\[
 \beta=\frac{66048RK}{UV(U+129)},\qquad
 \gamma=\frac{VTK}{RUY},\qquad
 \widetilde\rho=\frac{2064TK}{U(U+129)Y}.
\]
We verify the exact symbolic identities
\begin{equation}\label{comp:closure_identity}
 \begin{gathered}
 \mathcal B(\mathcal B(M_1,g_1,g_2,\beta),g_2,g_0,\gamma)
     =\widetilde\rho I,\\
 -\gamma\mathcal B(M_1,g_1,g_2,\beta)g_2
     =\widetilde\rho(g_0-g_2).
 \end{gathered}
\end{equation}
Structurally, in~\eqref{D:bfgs} the projection factors of
$M_2(\beta)=\mathcal B(M_1,g_1,g_2,\beta)$ are independent of $\beta$ while its
rank-one term is linear in $\beta$, so $M_2(\beta)=C+\beta E$. Requiring
$(C+\beta E)g_2$ to be parallel to $g_0-g_2$ is a single linear equation in
$\beta$ solved by the displayed formula; for a unit vector $w\perp(g_0-g_2)$,
rational evaluation gives $w^T(C+\beta E)w=\widetilde\rho$, and the displayed
$\gamma$ gives last step $s_2=-\gamma M_2(\beta)g_2=\widetilde\rho(g_0-g_2)$. At
the third update, $y_2=g_0-g_2$ is parallel to $s_2$, so the projection matrix
$I-s_2y_2^T/(y_2^Ts_2)$ is the orthogonal projector $\Pi=ww^T$ onto
$(g_0-g_2)^\perp$: its projected term equals $\Pi(C+\beta E)\Pi=\widetilde\rho\,\Pi$,
its rank-one term equals $s_2s_2^T/(y_2^Ts_2)=\widetilde\rho(I-\Pi)$, and their
sum is $\widetilde\rho I$, proving~\eqref{comp:closure_identity}.

At the root of~\eqref{comp:seed_polynomials}, $\beta=1/4$, $\gamma=1/16$, and
\[
 \widetilde\rho-\rho=
 \frac{2064R}{V(U+129)}\left(\gamma-\frac1{16}\right)=0.
\]
Direct substitution for $j=0,1$ and the second identity in~\eqref{comp:closure_identity}
for $j=2$ confirm that $s_j=-\alpha_jM_jg_j$ in~\eqref{D:steps} with
$(\alpha_0,\alpha_1,\alpha_2)=(1,1/4,1/16)$, while~\eqref{D:x0} solves
$s_0+s_1+s_2=-(1-\rho)z_0$, giving $z_3=\rho z_0$. Finally, the strict margins
certified below imply $(g_{j+1}-g_j)^Ts_j>0$ for $j=0,1,2$, so $M_0,M_1,M_2\succ0$
by \eqref{M:basic} in Chapter~\ref{comp:M}.
\end{proof}

\section{Rational interval certification of all nine queries}
Across one cycle, the unit-first search evaluates exactly nine trial points. We
bound them by exact rational interval arithmetic, using the standard rules
$[a,b]+[c,d]=[a+c,b+d]$, $-[a,b]=[-b,-a]$,
$[a,b][c,d]=[\min(ac,ad,bc,bd),\max(ac,ad,bc,bd)]$, and
$[a,b]^{-1}=[1/b,1/a]$ when $0\notin[a,b]$. Starting from the coarse rational
enclosure
\begin{equation}\label{comp:query_box}
 \begin{gathered}
 c\in\frac{[318481,318482]}{10^6},\qquad
 d\in\frac{[-1025233,-1025232]}{10^6},\\
 \rho\in\frac{[6019,6021]}{10^5}
 \end{gathered}
\end{equation}
implied by Lemma~\ref{comp:root_certificate} and~\eqref{D:rho}, we
evaluate $z_j$ from~\eqref{D:x0} and the search directions
\[
 p_0=(-1,0)^T,\qquad p_1=(127/256,1/16)^T,\qquad
 p_2=16\rho(1-c,-d)^T.
\]
For each row of the table in Section~\ref{comp:trial_table} with designated active branch $i$
and trial point $q=z_j+\alpha p_j$, we bound from below the active-branch,
slope, and acceptance/rejection margins
\[
 (g_i-g_k)^Tq\quad(k\ne i),\qquad -\alpha g_j^Tp_j,\qquad
 \begin{cases}
 g_j^Tz_j+\tfrac13\alpha g_j^Tp_j-g_i^Tq,&\text{acceptance},\\
 g_i^Tq-g_j^Tz_j,&\text{rejection},
 \end{cases}
\]
together with $\alpha g_i^Tp_j$ on accepted rows. Interval evaluation shows that
every margin in each row exceeds the corresponding entry of the positive
rational vector
\begin{equation}\label{comp:query_lower_bounds}
 \frac1{10^4}(729,\;2417,\;145,\;75,\;5964,\;2682,\;1041,\;221,\;20).
\end{equation}
For each trial value $G(z_j+\alpha p_j)$ in the table of Section~\ref{comp:trial_table} and each current value $G(z_j)$ in the display after it, the same interval operations give a rational enclosure of width at most $4\cdot10^{-4}$ (the twelve-digit decimals there are explanatory), as well as
positive active-branch gaps at $z_0,z_1,z_2$ (hence local differentiability at
every base point and query), and positivity of $3\pm g_i^Tq$, $2-(q)_2$, and
$g_i^Tq$ across all base points and queries, establishing~\eqref{D:trialbounds}
and completing the proof of the seed and query assertions in Chapter~\ref{comp:D}.

